\section{Proofs for the variance and twin-prime criteria}\label{app:criteria}

\subsection{The RH variance criterion}\label{app:rh-variance}

\begin{lemma}[The smoothed prime discrepancy]\label{cq:smoothed}
Fix a real \(m>1\), and set
\[
 H_m(v)=\sum_p(p+v)^{-m}
       -\int_2^\infty\frac{du}{(u+v)^m\log u},\qquad v>0.
\]
Put
\[
 B_m=\frac{\Gamma(1/2)\Gamma(m-1/2)}{2\Gamma(m)},\qquad
 S=1+\frac{\gamma_{\!E}}2-\frac12\log(4\pi).
\]
On RH, with nontrivial zeros counted with multiplicity,
\begin{equation}\label{cq:zeroseries}
 v^{m-1/2}\log v\,H_m(v)
 =-B_m-\sum_\rho
 \frac{\Gamma(\rho)\Gamma(m-\rho)}{\Gamma(m)}
 v^{i\Im\rho}+o(1).
\end{equation}
The series is absolutely convergent, uniformly in \(\log v\), and its absolute
value is at most \(2SB_m<B_m\). If RH fails, the expression on the left of
\eqref{cq:zeroseries} has limsup \(+\infty\) and liminf \(-\infty\).
\end{lemma}

\begin{proof}
We spell out the removal of the logarithmic prime weight. The classical
ingredients are the logarithmic derivative of zeta, its explicit formula,
the zero-counting bound and the RH estimate for the Chebyshev function; see
\cite[Chapters II, III, IX and XIV]{Titchmarsh1986}.
For \(a\geq0\), define
\[
 Q_{m,a}(v)=
 \sum_p\frac{(\log p)p^{-a}}{(p+v)^m}
 -\int_2^\infty\frac{u^{-a}}{(u+v)^m}\,du .
\]
Since \(\int_0^\infty u^{-a}\,da=1/\log u\), applying Tonelli separately to the
two nonnegative terms gives
\begin{equation}\label{cq:removeweight}
                      H_m(v)=\int_0^\infty Q_{m,a}(v)\,da.
\end{equation}
Assume RH and take \(0\leq a\leq a_0=1/8\). For \(1<c<m\), Mellin inversion is
\[
 \sum_{n\geq2}\frac{\Lambda(n)n^{-a}}{(n+v)^m}
 =\frac{1}{2\pi i}\int_{(c)}
 \frac{\Gamma(s)\Gamma(m-s)}{\Gamma(m)}
 \left(-\frac{\zeta'(s+a)}{\zeta(s+a)}\right)v^{s-m}\,ds.
\]
Move the line to \(\Re s=1/4\). The pole at \(1-a\) gives the continuous
integral over \((0,\infty)\); a zero \(\rho\) gives the residue
\[
 -\frac{\Gamma(\rho-a)\Gamma(m-\rho+a)}{\Gamma(m)}v^{\rho-a-m}.
\]
The remaining integral is \(O_m(v^{1/4-m})\), uniformly in \(a\).
Indeed, the new zeta arguments have real parts in \([1/4,3/8]\), separated
from all zeros on RH. Polynomial bounds for the logarithmic derivative and
exponential decay of the gamma factors control the vertical integral;
the horizontal segments may be chosen at the usual heights avoiding zeros.
Changing the continuous lower endpoint from \(0\) to \(2\) costs \(O_m(v^{-m})\).

We must remove the prime powers from the von Mangoldt sum. The squares
contribute
\[
 \sum_p\frac{(\log p)p^{-2a}}{(p^2+v)^m}
 =v^{1/2-a-m}
 \left\{\frac{\Gamma(1/2-a)\Gamma(m-1/2+a)}{2\Gamma(m)}+o_m(1)\right\}.
\]
This is the prime number theorem applied as a Stieltjes integral against
\(\vartheta(u)\), followed by \(u=\sqrt v\,y\). The rescaled kernels and
their derivatives have common integrable majorants for \(0\leq a\leq1/8\),
so the error is uniform. Powers of exponent at least three contribute
\(O_m(v^{1/3-m}\log^2 v)\), by Chebyshev's bound and dyadic summation.
Consequently
\[
 Q_{m,a}(v)=v^{1/2-a-m}
 \left\{-B_m(a)-\sum_\rho C_{m,\rho}(a)v^{i\Im\rho}+o_m(1)\right\},
\]
uniformly on the parameter interval, where
\[
 B_m(a)=\frac{\Gamma(1/2-a)\Gamma(m-1/2+a)}{2\Gamma(m)},\qquad
 C_{m,\rho}(a)=\frac{\Gamma(\rho-a)\Gamma(m-\rho+a)}{\Gamma(m)}.
\]
The zero sums of the coefficients and their first \(a\)-derivatives converge
uniformly. The elementary endpoint estimate
\[
 L\int_0^{a_0}e^{-aL}g(a)\,da=g(0)+o(1),\qquad L\longrightarrow\infty,
\]
therefore gives \eqref{cq:zeroseries} from this part of \eqref{cq:removeweight}.

The discarded parameter tail also needs a bound. With
\(E(u)=\vartheta(u)-u+2\), interpreted from \(2^-\), that tail equals
\[
 \int_{[2,\infty)}
 \frac{u^{-a_0}}{\log u\,(u+v)^m}\,dE(u).
\]
On RH, \(E(u)=O(\sqrt u\log^2u)\). Integration by parts bounds the last
display by
\[
 O_m\!\left(
 \int_2^\infty\frac{u^{-a_0-1/2}\log u}{(u+v)^m}\,du+
 \int_2^\infty\frac{u^{1/2-a_0}\log u}{(u+v)^{m+1}}\,du\right)
 =O_m(v^{1/2-a_0-m}\log v).
\]
This is \(o(v^{1/2-m}/\log v)\), as required.

Next we compare the square term with the zeros. The Hadamard product gives,
on RH,
\[
 \sum_{\gamma>0}\frac{1}{\gamma^2+1/4}=S=0.0230957\ldots .
\]
The gamma integral, reflection formula and
\(\cosh(\pi\gamma)\geq4(\gamma^2+1/4)^2\) imply
\[
 \sum_\rho\left|
 \frac{\Gamma(\rho)\Gamma(m-\rho)}{\Gamma(m)}
 \right|\leq 2B_m S .
\]
For the elementary hyperbolic bound it suffices to retain the constant,
quadratic and quartic terms of the series for \(\cosh\).
This proves the claimed strict negative margin on RH.

For the converse, initially in \(1<\Re s<m\), Mellin transformation gives
\begin{equation}\label{cq:mellinconverse}
 \int_0^\infty v^{m-s-1}H_m(v)\,dv
 =\frac{\Gamma(s)\Gamma(m-s)}{\Gamma(m)}F(s),
\end{equation}
where
\[
 F(s)=P(s)-J(s),\qquad
 P(s)=\sum_pp^{-s},\qquad
 J(s)=\int_2^\infty\frac{u^{-s}}{\log u}\,du .
\]
Write \(R(s)=\sum_{k\geq2}P(ks)/k\). This is analytic on \(\Re s>1/2\), and
continuation of the derivative of \(F\) gives
\begin{equation}\label{cq:Fderivative}
 F'(s)=\frac{\zeta'(s)}{\zeta(s)}-R'(s)+\frac{2^{1-s}}{s-1}.
\end{equation}
The singularities at the real point \(1\) cancel. There are no other real
singularities above \(1/2\). An off-line zero \(\rho\) with
\(\Re\rho>1/2\), however, gives a nonzero pole of \(F'\), hence a
nonremovable logarithmic singularity of \(F\).

Suppose the normalized discrepancy were eventually bounded below by \(-C\).
For a sufficiently large \(V>1\), the function
\[
 h(v)=v^m H_m(v)+C\,v^{1/2}/\log v
\]
would be nonnegative on \([V,\infty)\). Its Mellin integral
\(\int_V^\infty h(v)v^{-s-1}\,dv\) has an abscissa of convergence
\(\sigma_c\leq1\). The off-line singularity implies
\(\sigma_c\geq\Re\rho>1/2\): the added term is analytic to the right of
\(1/2\), and the short integral omitted from \eqref{cq:mellinconverse}
is analytic for \(\Re s<m\).

Landau's positivity argument forces a singularity at the real point
\(\sigma_c\), contradicting \eqref{cq:Fderivative}. Here is the argument
in the required form. If a Mellin transform of a nonnegative function
were analytic at its finite real abscissa, expand at a slightly larger
real point. The Taylor coefficients evaluated in the leftward direction
are integrals of nonnegative powers of \(\log v\). Tonelli then turns
convergence of the Taylor series just to the left of the abscissa into
convergence of the defining integral there, a contradiction.
Applying the same reasoning to \(-H_m\) rules out an eventual upper bound.
This proves both unbounded oscillations.
\end{proof}

Theorem~\ref{cq:rhvariance} follows by taking \(m=2\) and \(v=t-1\).
There is also a useful warning about the order of quantifiers. At any fixed
\(v>0\), the atom at \(2\) eventually dominates as \(m\) grows:
\[
 \sum_p(p+v)^{-m}\geq(v+2)^{-m},\qquad
 \int_2^\infty\frac{du}{(u+v)^m\log u}
 \leq\frac{(v+2)^{1-m}}{(m-1)\log2}.
\]
Thus \(H_m(v)>0\) whenever \((m-1)\log2>v+2\). Each fixed order can
eventually have the RH sign; one cannot require that sign for every order
on one common tail.



\subsection{Counting twin pairs with negative directions}\label{app:twin-index}

\begin{proof}[Proof of Theorem~\ref{cq:twins}]
The finite quotients are
\[
 F_y(t)=\prod_{2<p\leq y}(1-p^{-1})^3
                  \left(1+\frac3{t+p-1}\right).
\]
Each factor is completely monotone on \(t>-2\): its successive derivatives
alternate in sign. The products converge locally uniformly near that real
interval, because the terms of order \(1/p\) in their logarithms cancel.
Their derivatives converge too. Leibniz's rule proves complete monotonicity
of the limit; separating the factor for \(p=3\) makes all the inequalities
strict. This proves entrywise positivity.

The residue at the pole \(1-p\) is
\[
 b_p=\frac{K_o(4-p)}{K_o'(1-p)}
 =3(1-p^{-1})^3
  \prod_{\substack{q>2\\q\ne p}}
       (1-q^{-1})^3\left(1+\frac3{q-p}\right).
\]
For fixed \(p\), the tail factors are \(1+O_p(q^{-2})\).
Among the factors with \(q<p\), exactly one can be negative:
the factor with \(q=p-2\), if that prime exists. No factor vanishes,
since odd-prime differences cannot equal three. Hence
\begin{equation}\label{cq:twinsign}
                b_p<0\quad\Longleftrightarrow\quad p-2\text{ is prime}.
\end{equation}

To pass from residues to derivatives we need a uniform bound. Let \(b_{p,y}\)
be the finite-product residue. For \(q<p\), each factor in absolute value
is at most one. For \(p<q\leq2p\), comparison with all even distances gives
\[
 \prod_{p<q\leq2p}\left(1+\frac3{q-p}\right)
 \leq\prod_{1\leq j\leq p/2}\left(1+\frac3{2j}\right)\ll p^{3/2}.
\]
For \(q>2p\), the logarithm of the combined factor is at most
\[
             -\frac3q+\frac3{q-p}\leq\frac{6p}{q^2}.
\]
Even its sum over all integers is bounded. Thus, uniformly for \(y\geq p\),
\[
                      |b_{p,y}|\leq C p^{3/2}.
\]
Differentiate the finite partial fractions at least twice and use dominated
convergence:
\begin{equation}\label{cq:partialfractions}
 \frac{(-1)^k F^{(k)}(a)}{k!}
       =\sum_{p>2}\frac{b_p}{(a+p-1)^{k+1}},\qquad k\geq2.
\end{equation}
The absolute majorant for \(k=2\) is already bounded by
\(C\sum_{n\geq1}n^{-3/2}\). This explains why the matrix starts with the
second derivative.

Define the finite signed measure
\[
 \sigma_a=\sum_{p>2}\frac{b_p}{(a+p-1)^3}
                      \delta_{\,1/(a+p-1)} .
\]
Its support is compact, with its only possible accumulation point at zero.
Equation~\eqref{cq:partialfractions} identifies the matrix quadratic form as
\[
 c^{\mathsf T}\mathsf H_N(a)c
       =\int\left(\sum_{j=0}^{N-1}c_j u^j\right)^2\,d\sigma_a(u).
\]
If \(\sigma_a\) has \(J\) negative atoms, the negative part of this form
has rank at most \(J\), giving the upper bound on its negative index.

Conversely, choose any \(k\) negative atoms. They are isolated, so disjoint
continuous bump functions can select them and vanish at all other support
points. Normalize the functions so that their signed Gram matrix is
\(-I_k\). Uniform polynomial approximation preserves this strict negative
definiteness, because \(\sigma_a\) has finite total variation. The resulting
polynomials have some common finite degree and span a \(k\)-dimensional
negative subspace of a finite matrix. This proves the reverse bound,
whether there are finitely or infinitely many negative atoms.
Finally use \eqref{cq:twinsign} and the fact that successive matrices are
nested principal submatrices.
\end{proof}
